
Table \ref{tab:asQspiPer04a}, \ref{tab:asQspiPer04ab}, \ref{tab:asQspiPer04ac} shows the input and outputs of the QSPI kernel.
\begin{table}[H]
\caption{QSPI Kernel I/O (1)}
\label{tab:asQspiPer04a}
\centering
\begin{tabularx}{\textwidth}{|l |l |l |X|}
  \hline
  Pin & Direction & Width & Explanation \\
  \hline
  \hline
  clk\_i & in & 1 & Kernel clock \\
  \hline
  rst\_i & in & 1 & Kernel reset. Active high. \\
  \hline
\end{tabularx}
\end{table}

\begin{table}[H]
\caption{QSPI Kernel I/O (2)}
\label{tab:asQspiPer04ab}
\centering
\begin{tabularx}{\textwidth}{|l |l |l |X|}
  \hline
  Pin & Direction & Width & Explanation \\
  \hline
  \hline
  \multicolumn{4}{|l|}{Interface to BPI:} \\
  \hline
  ctrl\_reg11\_addr\_i & in & 1 & From QSPI-BPI (\asqspiCtrlRegExt register). Addr length (as member of an enum type) \\
  ctrl\_reg10\_cpol\_i & in & 1 & Clock polarity (enum type) \\
  ctrl\_reg9\_cpha\_i & in & 1 & Clock phase (enum type) \\
  ctrl\_reg8\_dual\_i & in & 1 & Dual data (enum type) \\
  ctrl\_reg7\_csh\_i & in & 1 & CS hold (enum type) \\
  ctrl\_reg6\_xip\_i & in & 1 & XIP (enum type) \\
  ctrl\_reg5\_ddr\_i & in & 1 & DDR (enum type) \\
  ctrl\_reg4\_quad\_i & in & 1 & Quad data (enum type) \\
  ctrl\_reg3\_dual\_i & in & 1 & Start (must be a pulse) (enum type) \\
  \hline
  cmd\_reg\_i(7:0) & in & 8 & From QSPI-BPI (\asqspiCmdRegExt register). Defines the SPI command opcode. \\
  \hline
  addr\_reg\_i(31:0) & in & 32 & From QSPI-BPI (\asqspiAddrRegExt register). Defines the target flash address. \\
  \hline
  len\_reg\_i(15:0) & in & 16 & From QSPI-BPI (\asqspiLenRegExt register). Defines the transfer length in bytes. \\
  \hline
  dummy\_reg\_i(5:0) & in & 6 & From QSPI-BPI (\asqspiDummyRegExt register). Defines the needed dummy cycles. \\
  \hline
  status\_reg\_o: & - & - & To \asqspiStatusRegExt reg. ...\\
  stat\_busy\_o & out & 1 & FSM active (level). \\
  stat\_done\_o & out & 1 & FSM in DONE state (pulse). \\
  stat\_error\_o & out & 1 & Error latched (level, until ICR). \\
  stat\_timeout\_o & out & 1 & Timeout (level, until ICR). \\
  tx\_empty\_i & in & 1 & From TX-FIFO. Empty indication. \\
  rx\_full\_i & in & 1 & From RX-FIFO. Full indication. \\
  tx\_rd\_o & out & 1 & Read strobe to TX-FIFO. \\
  rx\_wr\_o & out & 1 & Write strobe to RX-FIFO. \\
  timeout\_reg\_i(31:0) & in & 32 & Timeout level. \\
  \hline
  tx\_data\_i(63:0) & in & 64 & From QSPI-BPI (TX-FIFO). Data to be send. \\
  \hline
  rx\_data\_o(63:0) & out & 64 & To QSPI-BPI (RX-FIFO). Data received. \\
  \hline
\end{tabularx}
\end{table}

\begin{table}[H]
\caption{QSPI Kernel I/O (3)}
\label{tab:asQspiPer04ac}
\centering
\begin{tabularx}{\textwidth}{|l |l |l |X|}
  \hline
  Pin & Direction & Width & Explanation \\
  \hline
  \hline
  \multicolumn{4}{|l|}{Interface to PADs:} \\
  \hline
  \asqspiData& bidi & 4 &  To inout pins. QSPI data inout. Generated by the QSPI FSM. \\
  \hline
  \asqspiSCK & out & 1 &  To output pins. Clock for outside NOR-flash. Generated by the QSPI FSM.  \\
  \hline
  \asqspiCS & out & 1 &  To output pins. Chip select for outside NOR-flash. Generated by the QSPI FSM.  \\
  \hline
\end{tabularx}
\end{table}
